\phantomsection
\addcontentsline{toc}{chapter}{Rankings of Days}\label{ranking}
\fancyhead[C]{{\Large Ranking of Days}}
%\section{Ranking of Days}\label{ranking}
\noindent
\lett{F}{east} Days are ranked in order of priority. This is important, since it determines both what is prayed in the changeable parts of the Office (the Propers) and how to determine which Feast outranks the other(s) when Feast Days overlap on the same day.\par
The rankings for all of the Feast Days throughout the year are provided in the Kalendar (p. \pageref{kalendar}). For most people, it is expedient to order a yearly Ordo from the Antiochian Western Rite Vicariate which lays out the correct ordering for each day. However, to determine it without an Ordo, important rankings are described here.
\begin{itemize}
    \item Feast Days
        \begin{description}
            \item[First Class Double] (I Double)
            \item[Second Class Double] (II Double)
            \item[Greater Double]
            \item[Double] Feasts of Double and higher have I and II Evensong, beginning with I Evensong and ending with Compline of the following day.
            \item[Semidouble] Sundays and days within Common Octaves are Semidouble.\par
            The Semidouble Office begins with I Evensong and ends with Compline of the following day.
            \item[Memorial] No Office is to be said for a Memorial, but on the day on which they are noted in the Kalendar, Commemoration only is made of them at I Evensong and Mattins; except on II Doubles, on which Commemoration of the Memorial is omitted at Evensong; and on I Doubles, on which no Commemoration is made of Memorials.
            \item[Simple] The Office is Simple on weekdays when the Office is of the Feria, and on the Octave Day of a Simple Octave.\par
                The Simple Office begins with I Evensong, and ends with None of the following day.
        \end{description}
    \item Ferial Days
        \begin{description}
            \item[First Class Feria] No Feast Day may be celebrated. Begins at Mattins.
            %Not sure if this is true of the English Office:
                %The Office of the three Greater Ferias of Holy Week is said as directed in the Proper of the Season.
                \begin{itemize}
                    \item Ash Wednesday
                    \item Holy Monday
                    \item Holy Tuesday
                    \item Spy Wednesday
                \end{itemize}
            \item[Second Class Feria] Only I and II Doubles may be celebrated, lesser Feast Days being commemorated. Begins at Mattins.
            	\begin{itemize}
            		\item Ember Days
            		\item Rogation Monday
            		\item Weekdays in Lent \& Passion Week
            	\end{itemize}
            \item[Greater Feria] Begins at Mattins. 
                \begin{itemize}
                    \item Weekdays in Advent
                \end{itemize}
            \item[Feria] Begins when the Office of the preceding day ceases. The Ferial Office ends at None if a Double or Simple follow.
        \end{description}
    \item Octaves
        \begin{itemize}
            \item Privileged Octaves
                \begin{description}
                    \item[First Order Privileged Octaves] The Feast Day, and second \& third days of the Octave, are I Double. The rest of the Octave until the Octave Day is Semidouble. The Octave Day is I Double. Feast Days cannot outrank the Days within the Octave, nor the Octave Day.\par
                    A I or II Double is transferred to the first unhindered day after the Octave; other Feasts are commemorated, except on the Feast Day and the two following days.
                        \begin{itemize}
                            \item Easter Octave
                            \item Whitsun Octave
                        \end{itemize}
                    \item[Second Order Privileged Octaves] The Feast Day is I Double. The Days within the Octave are Semidouble, being only outranked by I Double, with the Octave commemorated. The Octave Day is Greater Double, and gives place only to a I Double of universal observance.
                        \begin{itemize}
                            \item Epiphany Octave
                            \item Corpus Christi Octave
                        \end{itemize}
                    \item[Third Order Privileged Octaves] Days within the Octave are trumped by any feast over Simple. Any Double Feast within this Octave has its Office with Commemoration of the Octave, but on the Octave Day of the Ascension only a I or II Double is kept, with Commemoration of the Octave.
                        \begin{itemize}
                            \item Christmas Octave
                            \item Ascension Octave
                        \end{itemize}
                \end{description}
            \item Common Octaves\par
            The Days within the Octave are Semidouble. A Double Feast occurring within a Common Octave has its Office with Commemoration of the Octave, unless otherwise noted.\par
            \textsc{Note,} Common Octaves are not commemorated during Lent.
                \begin{itemize}
                    \item Conception of the Blessed Virgin Mary Octave
                    \item Assumption of the Blessed Virgin Mary Octave
                    \item Nativity of St. John Baptist Octave
                    \item Solemnity of St. Joseph Octave
                    \item Sts. Peter \& Paul Octave
                    \item All Hallows Octave
                    \item The Octave of the principal patron saint of a church, cathedral, order, town, diocese, province, or nation. 
                \end{itemize}
            \item Simple Octaves\par
            The Feast Day is II Double. The Octave Day is Simple. The Days within the Octave are not commemorated.
                \begin{itemize}
                    \item St. Stephen Octave
                    \item St. John the Evangelist Octave
                    \item Holy Innocents Octave
                    \item St. Lawrence Octave
                    \item Nativity of the Blessed Virgin Mary Octave
                    \item Divine Compassion Octave
                    \item An Octave of Secondary Patrons
                \end{itemize}
        \end{itemize}
    \item Sundays
        \begin{description}
            \item[First Class] Sundays of the First Class give place to no Feast. A I or II Double Feast Day falling on a Sunday of the First Class is transferred to the first unhindered day.\par
            \textsc{Note,} When I or II Double Feast Day falls on (or is transferred to) Monday, I Evensong is of the Feast Day with a commemoration of the Sunday. Otherwise, it is of the Sunday with a commemoration of the Feast Day.
                \begin{itemize}
                    \item First Sunday of Advent
                    \item Sundays in Lent \& Passiontide
                    \item Easter Sunday
                    \item Low Sunday
                    \item Whitsunday
                \end{itemize}
            \item[Second Class] Sundays of the Second Class give place only to I Doubles, and the Sunday is commemorated.
                \begin{itemize}
                    \item Second Sunday of Advent
                    \item Third Sunday of Advent
                    \item Fourth Sunday of Advent
                    \item Septuagesima Sunday
                    \item Sexagesima Sunday
                    \item Quinquagesima Sunday
                \end{itemize}
            \item[All Other Sundays] On lesser Sundays, the Office is of the Sunday unless a Double I or II Class occur thereon, when the Office is of the Feast with Commemoration of the Sunday.
        \end{description}
    \item Privileged Vigils\par
        \begin{description}
            \item[First Class] Vigils of the First Class are preferred over any other Feast Day.
                \begin{itemize}
                    \item Christmas Even
                    \item Whitsun Even
                \end{itemize}
            \item[Second Class] Vigils of the Second Class are preferred over any Feast Day, except I \& II Doubles and Feasts of Our Lord.
                \begin{itemize}
                    \item Vigil of Epiphany
                \end{itemize}
        \end{description}
    \item Common Vigils\par
        Non-privileged Vigils are preferred over any Feast Day, except Doubles \& above or an Octave.
\end{itemize}

%\subsection{First \& Second Evensong}
%When I Evensong of a Feast Day overlaps with II Evensong of another Feast Day, one will always take priority.
%
%  	\begin{description}
 % 		\item[First Class Double with a Common Octave] I Evensong is always said.
 % 		\item[First Class Double] I Evensong is said, unless the II Evensong is of a First Class Double with a Common Octave. If the II Evensong Feast Day be of Greater Double or higher, Commemoration is made of it.
%  		\item[Second Class Double] I Evensong is said, unless the II Evensong be of a Second Class Double.\par
%  	\end{description}
